\section{Omitting types}
\label{sec:types-ultraproducts-03}
\label{sec:omitting-types}

\begin{definition}[Principal and nonprincipal types]
\label{def:principal-type}
Let \(p(\bar x)\) be a type consistent with a theory \(T\).  The type is
\emph{principal over \(T\)} iff some formula \(\theta(\bar x)\) satisfies
\begin{enumerate}
  \item \(T\cup\{\exists\bar x\,\theta(\bar x)\}\) is satisfiable; and
  \item for every \(\varphi(\bar x)\in p\),
        \(T\entails\forall\bar x(\theta\to\varphi)\).
\end{enumerate}
If no such formula exists, \(p\) is \emph{nonprincipal over \(T\)}.  This is
a property of a type relative to a theory, not a claim that a selected model
realizes or omits it.  It distinguishes a type forced by one consistent
formula from a merely finitely satisfiable collection, and that distinction is
the hypothesis of the theorem below.  When \(p\) is complete, it is principal
iff its point is isolated in \(S_n(T)\): an isolating formula
\(\theta\) gives \(\{p\}=[\theta]\), and a basic open singleton gives such a
formula.
\end{definition}

\begin{theorem}[Omitting types theorem]
\label{thm:omitting-types}
Let \(T\) be a satisfiable theory in a countable signature, and let
\(p_0,p_1,\ldots\) be a countable family of nonprincipal types consistent
with \(T\), each in a finite tuple of variables.  Then \(T\) has a countable
model that omits every \(p_n\).
\textit{Status: cited theorem.}
\end{theorem}

See \cite{Hodges1993,Marker2002} for complete proofs.  The Henkin construction
enumerates both witness requirements and pairs consisting of a term tuple and
a type \(p_n\).  At the latter requirement, nonprincipality permits adding the
negation of some formula in \(p_n\) at that tuple without destroying finite
satisfiability.  The term model therefore witnesses the theory while no tuple
realizes every formula of any listed type.

\begin{example}[A type that a countable model omits]
Take a signature with constants \(c_0,c_1,\ldots\), and let
\[
  T=\{c_m\ne c_n:m\ne n\},
  \qquad
  p(x)=\{x\ne c_n:n\in\N\}.
\]
The type is consistent and nonprincipal over \(T\).  The model whose domain
consists exactly of the distinct interpretations of the constants satisfies
\(T\) and omits \(p\), while a model with one additional unnamed element
realizes \(p\).  Thus omission is a property of a model, whereas
nonprincipality is the syntactic-semantic condition on the type used by
\cref{thm:omitting-types}.
\end{example}

